Para realizar una tabla se puede usar cualquiera de los entornos de \LaTeX\ diseñados a tal efecto como pueden ser \textbf{tabular}, \textbf{tabbing}, \textbf{longtable}, etc. Sin embargo estas tablas tienen que ser elementos flotantes con pie de tabla, referenciables y listables. Para ello se dispone del entorno \textbf{table}{[short]}\{label\}\{caption\}. El primer parámetro es opcional y por tanto, si aparece, debe ir entre corchetes y es texto corto que aparecerá en la lista de cuadros de texto; el segundo parámetro es una etiqueta para ser referenciada y no es opcional aunque puede dejarse en blanco; y el tercero es el texto que aparecerá como pie de tabla. Ejemplo de un cuadro de tabla \ref{TB:EJEMPLO}.

\begin{table}[Tabla de ejemplo]{TB:EJEMPLO}{Esta es una tabla de ejemplo en la que, internamente, se usa el entorno \textbf{tabular}.}
  \begin{tabular}{cccc}
    \hline
    \textbf{no} & \textbf{puedo} & \textbf{decir} & \textbf{nada} \\
    \hline \hline
    1.23 & 2.32 & 1.15 & 3.5 \\
    10.2 & 2.2 & 4.5 & 5.7 \\
    8.3 & 1.56 & 2.78 & 8.91 \\
    \hline
  \end{tabular}
\end{table}

La estética interna de las tablas es responsabilidad del autor pero se recomiendan diseños minimalistas y si se usan colores es aconsejable utilizar los colores complementarios definidos para la gama de colores correspondiente.

Así mismo, muchas veces es necesario introducir varias tablas juntas con sus correspondientes subpies de tabla. Para ello se puede utilizar el entorno \textbf{subtable} del paquete \textbf{subcaption} dentro del entorno \textbf{table}, igual que \textbf{subfigure} en las figuras: recibe la anchura de la subtabla y dentro se ponen la tabla, el pie con \textbf{{\textbackslash}caption} y la etiqueta con \textbf{{\textbackslash}label}. Un ejemplo con subtablas puede verse en la tabla \ref{TB:SUBTABLAS}.

\begin{table}[Tabla de ejemplo con subtablas]{TB:SUBTABLAS}{Esta es una tabla de ejemplo en la que se definen subtablas. Los pies de las subtablas no deben ser excesivos y se debe cargar toda la explicación posible en el pie de la tabla.}

    \begin{subtable}[t]{0.45\linewidth}
    \centering
    \begin{tabular}{cccc}
      \hline
      \textbf{no} & \textbf{puedo} & \textbf{decir} & \textbf{nada} \\
      \hline \hline
      1.23 & 2.32 & 1.15 & 3.5 \\
      10.2 & 2.2 & 4.5 & 5.7 \\
      8.3 & 1.56 & 2.78 & 8.91 \\
      \hline
    \end{tabular}
    \caption{Esta es una subtabla, y su pie no debe ser demasiado largo.}\label{STB:EJ1}
  \end{subtable}
    \begin{subtable}[t]{0.45\linewidth}
    \centering
    \begin{tabular}{cccc}
      \hline
      \textbf{no} & \textbf{puedo} & \textbf{decir} & \textbf{nada} \\
      \hline \hline
      11.39 & 1.21 & 5.15 & 2.9 \\
      5.2 & 4.8 & 9.43 & 1.7 \\
      7.3 & 6.35 & 0.11 & 3.13 \\
      \hline
    \end{tabular}
    \caption{Y este es otro ejemplo aunque con demasiado texto ambos.}\label{STB:EJ2}
  \end{subtable}
\end{table}
